\documentclass{report}
\usepackage{amsmath,amssymb}
\setlength{\parindent}{0mm}
\setlength{\parskip}{1em}
\begin{document}
\begin{center}
\rule{6in}{1pt} \
{\large Zhao-Li Shen \\
{\bf An efficient preconditioner based on a partial low-rank factorization for PageRank problems}}

Institute of Mathematics and Computing Science \\ University of Groningen \\ Nijenborgh 9 \\ P O Box 407 \\ 9700 AK Groningen \\ The Netherlands
\\
{\tt szlxiaoyao@163.com}\\
Bruno Carpentieri\end{center}

The PageRank problem used by Google for ranking web pages can be
mathematically stated as the following linear system
\begin{equation}\label{pagerank}
(I-\alpha P)x=v,
\end{equation}
where $0<\alpha<1$ is the damping factor, $P=\{p_{ij}\}$ is the
transition probability matrix representing the web structure where entry
$p_{ij}$ is the probability of transition from page $j$ to page $i$, $v$
is a personalization probability vector, and finally $x$ is the unknown
PageRank vector. With the rapid development of the Internet, the size of
system (\ref{pagerank}) is reaching order of millions or even billions
unknowns. For solving problems of this size, iterative methods are
mandatory to use. When the damping factor approaches 1, i.e., the model
reflects more closely the web structure, iterative solvers tend to suffer
from slow convergence due to the unfavourable spectral distribution and
the high condition number of system~(\ref{pagerank}) \cite{ref1,ref2}.
Thus preconditioning is necessary. Some preconditioning techniques that
have proved successful for general linear systems have been applied for
solving the PageRank problem. However, general purpose methods do not
exploit well any special property or structure of PageRank matrices.

In this study we exploit the similarity of pattern structures and the
special distribution of the numerical values on these patterns in the
transition matrix $P$ during the construction of preconditioner.
This similarity often appears due to the similar link distributions of
pages from same hosts. Note that the problem matrix $A=I-\alpha P$
maintains this similarity structure, except at the diagonal indexes. We
propose an economic compression algorithm that computes similarity
patterns in PageRank matrices and compresses matrix blocks of $A$ by a
low-rank matrix factorization, saving considerably storage and
computational cost. Accordingly, system matrix $A\in \mathbb{R}^{n\times
n}$ is reformulated as
\begin{equation}\label{low_rank}
A=D+FH,
\end{equation}
where $D\in \mathbb{R}^{n\times n}$ is more sparse than $A$, $F\in
\mathbb{N}^{n\times m}$ and $H \in \mathbb{R}^{m\times n}$ are low-rank
factors. We demonstrate that both the matrix $D$ and the matrix
$I+HD^{-1}F$ are unconditionally invertible for PageRank problems. Then
we propose a preconditioner $M$ based on the Woodbury formula of the form
\begin{equation}\label{preconditioner}
M\approx A^{-1}=D^{-1}(I-F(I+HD^{-1}F)^{-1}HD^{-1})
\end{equation}
that exploits effectively the computed structure (\ref{low_rank}).

Numerical experiments show that our algorithm can solve difficult
PageRank problems, i.e. for $\alpha=0.99,0.999$, with over one million
unknows efficiently, also in comparison to state-of-the-art solvers, and
may yield a high compression rate sometimes up to 40\%. In this talk we
present the theoretical background of the proposed method and we
illustrate numerical experiments of realistic PageRank simulations that
support our findings. Further improvement is under investigation.
{\small\begin{thebibliography}{99}
\bibitem{ref1}S.D. Kamvar and H. Haveliwala,
The Condition Number of the PageRank Problem, Technical
report, Stanford University, Stanford, CA, 2003.
\bibitem{ref2}G.M. Del Corso, A. Gulli, F. Romani, Comparison of Krylov
subspace methods on the
PageRank problem, \textit{J. Comput. Appl. Math.}, 210 (2007), pp. 159-166.
\end{thebibliography}}


\end{document}
